\cotefile{10}
\scanpage{42}
\mspage{\handwritten{10}}
\uline{pe de présentation finie, $H_j$ un $S_j$-préschéma de présentation finie à groupe d'\unsure{opérateurs|opératerus} $G_j$. Pour que $\mathbb{H}=H_j\handwritten{\times}_{S_j}S$ soit un $S$-préschéma en espaces homogènes sous $\deleted{\unsure{G|H}}\inserted{G=G_j\handwritten{\times}_{S_j}S}$ (resp. un fibré principal homogène sous $\unsure{G}\inserted{\deleted{\unsure{G=G_j\times_{S_j}S}}}$\unsure{)} pour la }\uuline{topologie fidèlement plate localement de présentation finie}\uline{, il faut et il suffit qu'il existe un indice $i\handwritten{\in}I$, $i\handwritten{\geq}j$, tel que $H_i=H_j\handwritten{\times}_{S_j}S_i$ soit un $S_i$-préschéma en espaces homogènes sous $\deleted{\unsure{G_i}}\inserted{G_i=G_j\handwritten{\times}_{S_j}S_i}$ (resp. un fibré principal homogène sous $G_i$).}

Compte tenu de (10.14) et de (EGA IV 8.8.2, 8.8.3 et 8.10.5) , \deleted{\unsure{l'exp|l'ex}} l'énoncé résulte de la propriété concernant les morphismes \deleted{\unsure{concernant|couvrant}}\inserted{couvrants} pour la topologie fppf rappelée en (10.0).

\figure{\begin{tikzpicture}[x=.5cm,y=.5cm]
\draw (0,11.6) .. controls (3.4,11.6) and (6.5,11.1) .. (5.1,10.2)
 .. controls (3.7,9.3) and (-1.8,7.2) .. (-.1,7.1)
 .. controls (1.9,7.0) and (4.1,7.6) .. (3.7,5.8)
 .. controls (3.5,4.8) and (2.1,3.8) .. (2.7,2.9)
 .. controls (3.4,1.9) and (5.0,1.1) .. (4.5,0);
\draw (2.45,11.7) .. controls (2.5,7.2) and (2.9,4.0) .. (3.75,.1);
\end{tikzpicture}}
% archive sheet number: 41
